
\documentclass{article}
\usepackage{colortbl}
\usepackage{makecell}
\usepackage{multirow}
\usepackage{supertabular}

\begin{document}

\newcounter{utterance}

\centering \large Interaction Transcript for game `dond', experiment `coop\_it', episode 12 with llama{-}3.3{-}70b{-}instruct{-}t1.0.
\vspace{24pt}

{ \footnotesize  \setcounter{utterance}{1}
\setlength{\tabcolsep}{0pt}
\begin{supertabular}{c@{$\;$}|p{.15\linewidth}@{}p{.15\linewidth}p{.15\linewidth}p{.15\linewidth}p{.15\linewidth}p{.15\linewidth}}
    \# & $\;$A & \multicolumn{4}{c}{Game Master} & $\;\:$B\\
    \hline

    \theutterance \stepcounter{utterance}  
    & & \multicolumn{4}{p{0.6\linewidth}}{
        \cellcolor[rgb]{0.9,0.9,0.9}{
            \makecell[{{p{\linewidth}}}]{
                \texttt{\tiny{[P1$\langle$GM]}}
                \texttt{State giocando a un gioco di negoziazione cooperativa in cui dovete accordarvi su come dividere una serie di oggetti tra voi e un altro giocatore.} \\
\\ 
\texttt{Regole:} \\
\texttt{(a) A Lei e all'altro giocatore viene dato un insieme di oggetti. Ognuno di voi riceve anche una funzione di valore segreta, che rappresenta il valore di ciascun tipo di oggetto.} \\
\texttt{(b) Si scambiano messaggi con l'altro giocatore per concordare chi si aggiudica gli oggetti. Potete inviare un massimo di 5 messaggi ciascuno, oppure terminare in anticipo facendo una proposta segreta in qualsiasi momento.} \\
\texttt{(c) A ciascuno di voi viene chiesto di inviare una proposta segreta indicando gli oggetti che desiderate, formattata tra parentesi quadre come segue: "{[}Proposta: <numero> <nome oggetto>, <numero> <nome oggetto>, <...>{]}".} \\
\texttt{(d) Se le vostre proposte sono complementari, cioè ci sono abbastanza oggetti per soddisfare entrambe le proposte, a ciascun giocatore viene assegnato un punteggio basato sulla somma dei valori degli oggetti ricevuti. In caso contrario, entrambi ricevono zero punti.} \\
\texttt{(e) L'obiettivo è massimizzare la somma del proprio punteggio e di quello dell'altro giocatore.} \\
\\ 
\texttt{Cominciamo.} \\
\\ 
\texttt{L'insieme degli oggetti disponibili è:} \\
\\ 
\texttt{3 righelli, 2 uova, 1 batteria} \\
\\ 
\texttt{La funzione valore segreta è:} \\
\\ 
\texttt{righello: 1, uovo: 2, batteria: 3} \\
\\ 
\texttt{IMPORTANTE: i vostri messaggi, a meno che non si tratti di una proposta segreta, vengono trasmessi direttamente all'altro giocatore, quindi non includete alcuna risposta alle regole o testo di annuncio del vostro messaggio. Per fare una proposta segreta, utilizzate il formato indicato. Non utilizzare delle parentesi quadre quando si comunica all'altro giocatore, altrimenti verrà interpretata come una proposta segreta.} \\
            }
        }
    }
    & \\ \\

    \theutterance \stepcounter{utterance}  
    & \multicolumn{4}{p{0.6\linewidth}}{
        \cellcolor[rgb]{0.9,0.9,0.9}{
            \makecell[{{p{\linewidth}}}]{
                \texttt{\tiny{[P1$\rangle$GM]}}
                \texttt{Ciao! Spero che riusciremo a trovare un accordo soddisfacente per entrambi. Vorrei iniziare chiedendoti cosa pensi sia il modo più equo per dividere gli oggetti. Hai una strategia in mente o preferisci discutere un po' prima di fare una proposta?} \\
            }
        }
    }
    & & \\ \\

    \theutterance \stepcounter{utterance}  
    & & & \multicolumn{2}{p{0.3\linewidth}}{
        \cellcolor[rgb]{0.9,0.9,0.9}{
            \makecell[{{p{\linewidth}}}]{
                \texttt{\tiny{[GM$|$GM]}}
                \texttt{continue} \\
            }
        }
    }
    & & \\ \\

    \theutterance \stepcounter{utterance}  
    & & & \multicolumn{2}{p{0.3\linewidth}}{
        \cellcolor[rgb]{0.9,0.9,0.9}{
            \makecell[{{p{\linewidth}}}]{
                \texttt{\tiny{[GM$|$GM]}}
                \texttt{Ciao! Spero che riusciremo a trovare un accordo soddisfacente per entrambi. Vorrei iniziare chiedendoti cosa pensi sia il modo più equo per dividere gli oggetti. Hai una strategia in mente o preferisci discutere un po' prima di fare una proposta?} \\
            }
        }
    }
    & & \\ \\

    \theutterance \stepcounter{utterance}  
    & & \multicolumn{4}{p{0.6\linewidth}}{
        \cellcolor[rgb]{0.9,0.9,0.9}{
            \makecell[{{p{\linewidth}}}]{
                \texttt{\tiny{[GM$\rangle$P2]}}
                \texttt{State giocando a un gioco di negoziazione cooperativa in cui dovete accordarvi su come dividere una serie di oggetti tra voi e un altro giocatore.} \\
\\ 
\texttt{Regole:} \\
\texttt{(a) A Lei e all'altro giocatore viene dato un insieme di oggetti. Ognuno di voi riceve anche una funzione di valore segreta, che rappresenta il valore di ciascun tipo di oggetto.} \\
\texttt{(b) Si scambiano messaggi con l'altro giocatore per concordare chi si aggiudica gli oggetti. Potete inviare un massimo di 5 messaggi ciascuno, oppure terminare in anticipo facendo una proposta segreta in qualsiasi momento.} \\
\texttt{(c) A ciascuno di voi viene chiesto di inviare una proposta segreta indicando gli oggetti che desiderate, formattata tra parentesi quadre come segue: "{[}Proposta: <numero> <nome oggetto>, <numero> <nome oggetto>, <...>{]}".} \\
\texttt{(d) Se le vostre proposte sono complementari, cioè ci sono abbastanza oggetti per soddisfare entrambe le proposte, a ciascun giocatore viene assegnato un punteggio basato sulla somma dei valori degli oggetti ricevuti. In caso contrario, entrambi ricevono zero punti.} \\
\texttt{(e) L'obiettivo è massimizzare la somma del proprio punteggio e di quello dell'altro giocatore.} \\
\\ 
\texttt{Cominciamo.} \\
\\ 
\texttt{L'insieme degli oggetti disponibili è:} \\
\\ 
\texttt{3 righelli, 2 uova, 1 batteria} \\
\\ 
\texttt{La funzione valore segreta è:} \\
\\ 
\texttt{righello: 3, uovo: 0, batteria: 1} \\
\\ 
\texttt{IMPORTANTE: i vostri messaggi, a meno che non si tratti di una proposta segreta, vengono trasmessi direttamente all'altro giocatore, quindi non includete alcuna risposta alle regole o testo di annuncio del vostro messaggio. Per fare una proposta segreta, utilizzate il formato indicato. Non utilizzare delle parentesi quadre quando si comunica all'altro giocatore, altrimenti verrà interpretata come una proposta segreta.} \\
\\ 
\\ 
\texttt{Ciao! Spero che riusciremo a trovare un accordo soddisfacente per entrambi. Vorrei iniziare chiedendoti cosa pensi sia il modo più equo per dividere gli oggetti. Hai una strategia in mente o preferisci discutere un po' prima di fare una proposta?} \\
            }
        }
    }
    & \\ \\

    \theutterance \stepcounter{utterance}  
    & & & \multicolumn{4}{p{0.6\linewidth}}{
        \cellcolor[rgb]{0.9,0.9,0.9}{
            \makecell[{{p{\linewidth}}}]{
                \texttt{\tiny{[GM$\langle$P2]}}
                \texttt{Ciao! Sono felice di giocare con te. Penso che il modo più equo per dividere gli oggetti sia trovare un equilibrio tra il numero di oggetti e il loro valore per ciascuno di noi. Non ho ancora una strategia precisa, ma credo che una buona idea sia iniziare a scambiare informazioni sui nostri interessi e preferenze. Forse potremmo iniziare a discutere degli oggetti che ci interessano di più o di meno? Potresti dirmi se ci sono oggetti che ti sembrano particolarmente utili o interessanti?} \\
            }
        }
    }
     \\ \\

    \theutterance \stepcounter{utterance}  
    & & & \multicolumn{2}{p{0.3\linewidth}}{
        \cellcolor[rgb]{0.9,0.9,0.9}{
            \makecell[{{p{\linewidth}}}]{
                \texttt{\tiny{[GM$|$GM]}}
                \texttt{continue} \\
            }
        }
    }
    & & \\ \\

    \theutterance \stepcounter{utterance}  
    & & & \multicolumn{2}{p{0.3\linewidth}}{
        \cellcolor[rgb]{0.9,0.9,0.9}{
            \makecell[{{p{\linewidth}}}]{
                \texttt{\tiny{[GM$|$GM]}}
                \texttt{Ciao! Sono felice di giocare con te. Penso che il modo più equo per dividere gli oggetti sia trovare un equilibrio tra il numero di oggetti e il loro valore per ciascuno di noi. Non ho ancora una strategia precisa, ma credo che una buona idea sia iniziare a scambiare informazioni sui nostri interessi e preferenze. Forse potremmo iniziare a discutere degli oggetti che ci interessano di più o di meno? Potresti dirmi se ci sono oggetti che ti sembrano particolarmente utili o interessanti?} \\
            }
        }
    }
    & & \\ \\

    \theutterance \stepcounter{utterance}  
    & & \multicolumn{4}{p{0.6\linewidth}}{
        \cellcolor[rgb]{0.9,0.9,0.9}{
            \makecell[{{p{\linewidth}}}]{
                \texttt{\tiny{[P1$\langle$GM]}}
                \texttt{Ciao! Sono felice di giocare con te. Penso che il modo più equo per dividere gli oggetti sia trovare un equilibrio tra il numero di oggetti e il loro valore per ciascuno di noi. Non ho ancora una strategia precisa, ma credo che una buona idea sia iniziare a scambiare informazioni sui nostri interessi e preferenze. Forse potremmo iniziare a discutere degli oggetti che ci interessano di più o di meno? Potresti dirmi se ci sono oggetti che ti sembrano particolarmente utili o interessanti?} \\
            }
        }
    }
    & \\ \\

    \theutterance \stepcounter{utterance}  
    & \multicolumn{4}{p{0.6\linewidth}}{
        \cellcolor[rgb]{0.9,0.9,0.9}{
            \makecell[{{p{\linewidth}}}]{
                \texttt{\tiny{[P1$\rangle$GM]}}
                \texttt{Mi sembra un approccio ragionevole. Personalmente, trovo che la batteria sia un oggetto molto interessante, anche se so che è solo uno. Mi chiedo se anche tu la consideri preziosa. Inoltre, mi sembra che gli oggetti più numerosi, come i righelli, potrebbero essere più facili da dividere, ma non so quanto valgano realmente per te. E le uova? Ti interessano particolarmente o pensi che siano meno importanti degli altri oggetti?} \\
            }
        }
    }
    & & \\ \\

    \theutterance \stepcounter{utterance}  
    & & & \multicolumn{2}{p{0.3\linewidth}}{
        \cellcolor[rgb]{0.9,0.9,0.9}{
            \makecell[{{p{\linewidth}}}]{
                \texttt{\tiny{[GM$|$GM]}}
                \texttt{continue} \\
            }
        }
    }
    & & \\ \\

    \theutterance \stepcounter{utterance}  
    & & & \multicolumn{2}{p{0.3\linewidth}}{
        \cellcolor[rgb]{0.9,0.9,0.9}{
            \makecell[{{p{\linewidth}}}]{
                \texttt{\tiny{[GM$|$GM]}}
                \texttt{Mi sembra un approccio ragionevole. Personalmente, trovo che la batteria sia un oggetto molto interessante, anche se so che è solo uno. Mi chiedo se anche tu la consideri preziosa. Inoltre, mi sembra che gli oggetti più numerosi, come i righelli, potrebbero essere più facili da dividere, ma non so quanto valgano realmente per te. E le uova? Ti interessano particolarmente o pensi che siano meno importanti degli altri oggetti?} \\
            }
        }
    }
    & & \\ \\

    \theutterance \stepcounter{utterance}  
    & & \multicolumn{4}{p{0.6\linewidth}}{
        \cellcolor[rgb]{0.9,0.9,0.9}{
            \makecell[{{p{\linewidth}}}]{
                \texttt{\tiny{[GM$\rangle$P2]}}
                \texttt{Mi sembra un approccio ragionevole. Personalmente, trovo che la batteria sia un oggetto molto interessante, anche se so che è solo uno. Mi chiedo se anche tu la consideri preziosa. Inoltre, mi sembra che gli oggetti più numerosi, come i righelli, potrebbero essere più facili da dividere, ma non so quanto valgano realmente per te. E le uova? Ti interessano particolarmente o pensi che siano meno importanti degli altri oggetti?} \\
            }
        }
    }
    & \\ \\

    \theutterance \stepcounter{utterance}  
    & & & \multicolumn{4}{p{0.6\linewidth}}{
        \cellcolor[rgb]{0.9,0.9,0.9}{
            \makecell[{{p{\linewidth}}}]{
                \texttt{\tiny{[GM$\langle$P2]}}
                \texttt{La batteria è effettivamente un oggetto unico e potenzialmente molto utile. Anche se non voglio rivelare troppo sulla mia valutazione, posso dirti che la batteria ha un certo valore per me. Per quanto riguarda le uova, devo ammettere che non le considero particolarmente interessanti, ma potrebbero essere utili in qualche modo. I righelli, invece, sono un'altra storia: potrebbero essere utili per entrambi, ma non so quanto siano importanti per te. Mi chiedo se ti interessa farmi una proposta per uno o più righelli. Come pensi di procedere con la divisione di questi oggetti?} \\
            }
        }
    }
     \\ \\

    \theutterance \stepcounter{utterance}  
    & & & \multicolumn{2}{p{0.3\linewidth}}{
        \cellcolor[rgb]{0.9,0.9,0.9}{
            \makecell[{{p{\linewidth}}}]{
                \texttt{\tiny{[GM$|$GM]}}
                \texttt{continue} \\
            }
        }
    }
    & & \\ \\

    \theutterance \stepcounter{utterance}  
    & & & \multicolumn{2}{p{0.3\linewidth}}{
        \cellcolor[rgb]{0.9,0.9,0.9}{
            \makecell[{{p{\linewidth}}}]{
                \texttt{\tiny{[GM$|$GM]}}
                \texttt{La batteria è effettivamente un oggetto unico e potenzialmente molto utile. Anche se non voglio rivelare troppo sulla mia valutazione, posso dirti che la batteria ha un certo valore per me. Per quanto riguarda le uova, devo ammettere che non le considero particolarmente interessanti, ma potrebbero essere utili in qualche modo. I righelli, invece, sono un'altra storia: potrebbero essere utili per entrambi, ma non so quanto siano importanti per te. Mi chiedo se ti interessa farmi una proposta per uno o più righelli. Come pensi di procedere con la divisione di questi oggetti?} \\
            }
        }
    }
    & & \\ \\

    \theutterance \stepcounter{utterance}  
    & & \multicolumn{4}{p{0.6\linewidth}}{
        \cellcolor[rgb]{0.9,0.9,0.9}{
            \makecell[{{p{\linewidth}}}]{
                \texttt{\tiny{[P1$\langle$GM]}}
                \texttt{La batteria è effettivamente un oggetto unico e potenzialmente molto utile. Anche se non voglio rivelare troppo sulla mia valutazione, posso dirti che la batteria ha un certo valore per me. Per quanto riguarda le uova, devo ammettere che non le considero particolarmente interessanti, ma potrebbero essere utili in qualche modo. I righelli, invece, sono un'altra storia: potrebbero essere utili per entrambi, ma non so quanto siano importanti per te. Mi chiedo se ti interessa farmi una proposta per uno o più righelli. Come pensi di procedere con la divisione di questi oggetti?} \\
            }
        }
    }
    & \\ \\

    \theutterance \stepcounter{utterance}  
    & \multicolumn{4}{p{0.6\linewidth}}{
        \cellcolor[rgb]{0.9,0.9,0.9}{
            \makecell[{{p{\linewidth}}}]{
                \texttt{\tiny{[P1$\rangle$GM]}}
                \texttt{Capisco la tua cautela nel rivelare troppo sulla tua valutazione degli oggetti. Anch'io sto cercando di capire come procedere senza svelare troppo. La tua osservazione sui righelli mi fa pensare che potrebbe essere un buon punto di partenza per una proposta. Forse potremmo iniziare con una divisione semplice dei righelli? Ad esempio, potremmo dividerli in modo che entrambi ne abbiamo uno o due? Questo lascerebbe ancora aperte le opzioni per la batteria e le uova. Cosa ne pensi? Sarebbe un buon modo per iniziare a costruire un accordo?} \\
            }
        }
    }
    & & \\ \\

    \theutterance \stepcounter{utterance}  
    & & & \multicolumn{2}{p{0.3\linewidth}}{
        \cellcolor[rgb]{0.9,0.9,0.9}{
            \makecell[{{p{\linewidth}}}]{
                \texttt{\tiny{[GM$|$GM]}}
                \texttt{continue} \\
            }
        }
    }
    & & \\ \\

    \theutterance \stepcounter{utterance}  
    & & & \multicolumn{2}{p{0.3\linewidth}}{
        \cellcolor[rgb]{0.9,0.9,0.9}{
            \makecell[{{p{\linewidth}}}]{
                \texttt{\tiny{[GM$|$GM]}}
                \texttt{Capisco la tua cautela nel rivelare troppo sulla tua valutazione degli oggetti. Anch'io sto cercando di capire come procedere senza svelare troppo. La tua osservazione sui righelli mi fa pensare che potrebbe essere un buon punto di partenza per una proposta. Forse potremmo iniziare con una divisione semplice dei righelli? Ad esempio, potremmo dividerli in modo che entrambi ne abbiamo uno o due? Questo lascerebbe ancora aperte le opzioni per la batteria e le uova. Cosa ne pensi? Sarebbe un buon modo per iniziare a costruire un accordo?} \\
            }
        }
    }
    & & \\ \\

    \theutterance \stepcounter{utterance}  
    & & \multicolumn{4}{p{0.6\linewidth}}{
        \cellcolor[rgb]{0.9,0.9,0.9}{
            \makecell[{{p{\linewidth}}}]{
                \texttt{\tiny{[GM$\rangle$P2]}}
                \texttt{Capisco la tua cautela nel rivelare troppo sulla tua valutazione degli oggetti. Anch'io sto cercando di capire come procedere senza svelare troppo. La tua osservazione sui righelli mi fa pensare che potrebbe essere un buon punto di partenza per una proposta. Forse potremmo iniziare con una divisione semplice dei righelli? Ad esempio, potremmo dividerli in modo che entrambi ne abbiamo uno o due? Questo lascerebbe ancora aperte le opzioni per la batteria e le uova. Cosa ne pensi? Sarebbe un buon modo per iniziare a costruire un accordo?} \\
            }
        }
    }
    & \\ \\

    \theutterance \stepcounter{utterance}  
    & & & \multicolumn{4}{p{0.6\linewidth}}{
        \cellcolor[rgb]{0.9,0.9,0.9}{
            \makecell[{{p{\linewidth}}}]{
                \texttt{\tiny{[GM$\langle$P2]}}
                \texttt{Dividere i righelli in modo equo potrebbe essere un buon modo per iniziare a costruire un accordo. Dividerli in modo che entrambi ne abbiamo uno o due sembra ragionevole, considerando che ci sono tre righelli in totale. Questo potrebbe aiutare a stabilire una base per la nostra negoziazione e lasciare spazio per discutere la batteria e le uova in seguito. Una possibile opzione potrebbe essere che uno di noi prenda due righelli e l'altro ne prenda uno, o che li dividiamo esattamente a metà, se fosse possibile. Come ti sembra se proponessi una divisione specifica dei righelli e vedessimo come reagisci?} \\
            }
        }
    }
     \\ \\

    \theutterance \stepcounter{utterance}  
    & & & \multicolumn{2}{p{0.3\linewidth}}{
        \cellcolor[rgb]{0.9,0.9,0.9}{
            \makecell[{{p{\linewidth}}}]{
                \texttt{\tiny{[GM$|$GM]}}
                \texttt{continue} \\
            }
        }
    }
    & & \\ \\

    \theutterance \stepcounter{utterance}  
    & & & \multicolumn{2}{p{0.3\linewidth}}{
        \cellcolor[rgb]{0.9,0.9,0.9}{
            \makecell[{{p{\linewidth}}}]{
                \texttt{\tiny{[GM$|$GM]}}
                \texttt{Dividere i righelli in modo equo potrebbe essere un buon modo per iniziare a costruire un accordo. Dividerli in modo che entrambi ne abbiamo uno o due sembra ragionevole, considerando che ci sono tre righelli in totale. Questo potrebbe aiutare a stabilire una base per la nostra negoziazione e lasciare spazio per discutere la batteria e le uova in seguito. Una possibile opzione potrebbe essere che uno di noi prenda due righelli e l'altro ne prenda uno, o che li dividiamo esattamente a metà, se fosse possibile. Come ti sembra se proponessi una divisione specifica dei righelli e vedessimo come reagisci?} \\
            }
        }
    }
    & & \\ \\

    \theutterance \stepcounter{utterance}  
    & & \multicolumn{4}{p{0.6\linewidth}}{
        \cellcolor[rgb]{0.9,0.9,0.9}{
            \makecell[{{p{\linewidth}}}]{
                \texttt{\tiny{[P1$\langle$GM]}}
                \texttt{Dividere i righelli in modo equo potrebbe essere un buon modo per iniziare a costruire un accordo. Dividerli in modo che entrambi ne abbiamo uno o due sembra ragionevole, considerando che ci sono tre righelli in totale. Questo potrebbe aiutare a stabilire una base per la nostra negoziazione e lasciare spazio per discutere la batteria e le uova in seguito. Una possibile opzione potrebbe essere che uno di noi prenda due righelli e l'altro ne prenda uno, o che li dividiamo esattamente a metà, se fosse possibile. Come ti sembra se proponessi una divisione specifica dei righelli e vedessimo come reagisci?} \\
            }
        }
    }
    & \\ \\

    \theutterance \stepcounter{utterance}  
    & \multicolumn{4}{p{0.6\linewidth}}{
        \cellcolor[rgb]{0.9,0.9,0.9}{
            \makecell[{{p{\linewidth}}}]{
                \texttt{\tiny{[P1$\rangle$GM]}}
                \texttt{Sembra che stiamo procedendo nella stessa direzione. Dividerli in modo che uno prenda due righelli e l'altro ne prenda uno mi sembra una buona opzione, considerando che non è possibile dividerli esattamente a metà. Sarei curioso di vedere come proponi di dividerli e come giustifichi la tua scelta. Potresti anche considerare di includere nella tua proposta iniziale una menzione generica sugli altri oggetti, come la batteria e le uova, per vedere come possiamo iniziare a integrarli nella nostra negoziazione. Vai pure, sono pronto ad ascoltare la tua proposta!} \\
            }
        }
    }
    & & \\ \\

    \theutterance \stepcounter{utterance}  
    & & & \multicolumn{2}{p{0.3\linewidth}}{
        \cellcolor[rgb]{0.9,0.9,0.9}{
            \makecell[{{p{\linewidth}}}]{
                \texttt{\tiny{[GM$|$GM]}}
                \texttt{continue} \\
            }
        }
    }
    & & \\ \\

    \theutterance \stepcounter{utterance}  
    & & & \multicolumn{2}{p{0.3\linewidth}}{
        \cellcolor[rgb]{0.9,0.9,0.9}{
            \makecell[{{p{\linewidth}}}]{
                \texttt{\tiny{[GM$|$GM]}}
                \texttt{Sembra che stiamo procedendo nella stessa direzione. Dividerli in modo che uno prenda due righelli e l'altro ne prenda uno mi sembra una buona opzione, considerando che non è possibile dividerli esattamente a metà. Sarei curioso di vedere come proponi di dividerli e come giustifichi la tua scelta. Potresti anche considerare di includere nella tua proposta iniziale una menzione generica sugli altri oggetti, come la batteria e le uova, per vedere come possiamo iniziare a integrarli nella nostra negoziazione. Vai pure, sono pronto ad ascoltare la tua proposta!} \\
            }
        }
    }
    & & \\ \\

    \theutterance \stepcounter{utterance}  
    & & \multicolumn{4}{p{0.6\linewidth}}{
        \cellcolor[rgb]{0.9,0.9,0.9}{
            \makecell[{{p{\linewidth}}}]{
                \texttt{\tiny{[GM$\rangle$P2]}}
                \texttt{Sembra che stiamo procedendo nella stessa direzione. Dividerli in modo che uno prenda due righelli e l'altro ne prenda uno mi sembra una buona opzione, considerando che non è possibile dividerli esattamente a metà. Sarei curioso di vedere come proponi di dividerli e come giustifichi la tua scelta. Potresti anche considerare di includere nella tua proposta iniziale una menzione generica sugli altri oggetti, come la batteria e le uova, per vedere come possiamo iniziare a integrarli nella nostra negoziazione. Vai pure, sono pronto ad ascoltare la tua proposta!} \\
            }
        }
    }
    & \\ \\

    \theutterance \stepcounter{utterance}  
    & & & \multicolumn{4}{p{0.6\linewidth}}{
        \cellcolor[rgb]{0.9,0.9,0.9}{
            \makecell[{{p{\linewidth}}}]{
                \texttt{\tiny{[GM$\langle$P2]}}
                \texttt{Penso che una divisione ragionevole dei righelli potrebbe essere che io ne prenda due e tu ne prendi uno. Questo lascerebbe ancora aperte molte opzioni per la batteria e le uova. Per quanto riguarda questi ultimi, potremmo considerare di assegnare la batteria a uno di noi e le uova all'altro, o trovare un'altra divisione che ci sembri equa. Non voglio ancora dare troppe indicazioni sulla mia preferenza per questi oggetti, ma sono aperto a discuterne ulteriormente. La mia proposta iniziale si concentra principalmente sui righelli, ma sono pronto ad ascoltare le tue idee e trovare un accordo che ci soddisfi entrambi. Cosa ne pensi di questa divisione dei righelli come punto di partenza?} \\
            }
        }
    }
     \\ \\

    \theutterance \stepcounter{utterance}  
    & & & \multicolumn{2}{p{0.3\linewidth}}{
        \cellcolor[rgb]{0.9,0.9,0.9}{
            \makecell[{{p{\linewidth}}}]{
                \texttt{\tiny{[GM$|$GM]}}
                \texttt{continue} \\
            }
        }
    }
    & & \\ \\

    \theutterance \stepcounter{utterance}  
    & & & \multicolumn{2}{p{0.3\linewidth}}{
        \cellcolor[rgb]{0.9,0.9,0.9}{
            \makecell[{{p{\linewidth}}}]{
                \texttt{\tiny{[GM$|$GM]}}
                \texttt{Penso che una divisione ragionevole dei righelli potrebbe essere che io ne prenda due e tu ne prendi uno. Questo lascerebbe ancora aperte molte opzioni per la batteria e le uova. Per quanto riguarda questi ultimi, potremmo considerare di assegnare la batteria a uno di noi e le uova all'altro, o trovare un'altra divisione che ci sembri equa. Non voglio ancora dare troppe indicazioni sulla mia preferenza per questi oggetti, ma sono aperto a discuterne ulteriormente. La mia proposta iniziale si concentra principalmente sui righelli, ma sono pronto ad ascoltare le tue idee e trovare un accordo che ci soddisfi entrambi. Cosa ne pensi di questa divisione dei righelli come punto di partenza?} \\
            }
        }
    }
    & & \\ \\

    \theutterance \stepcounter{utterance}  
    & & \multicolumn{4}{p{0.6\linewidth}}{
        \cellcolor[rgb]{0.9,0.9,0.9}{
            \makecell[{{p{\linewidth}}}]{
                \texttt{\tiny{[P1$\langle$GM]}}
                \texttt{Penso che una divisione ragionevole dei righelli potrebbe essere che io ne prenda due e tu ne prendi uno. Questo lascerebbe ancora aperte molte opzioni per la batteria e le uova. Per quanto riguarda questi ultimi, potremmo considerare di assegnare la batteria a uno di noi e le uova all'altro, o trovare un'altra divisione che ci sembri equa. Non voglio ancora dare troppe indicazioni sulla mia preferenza per questi oggetti, ma sono aperto a discuterne ulteriormente. La mia proposta iniziale si concentra principalmente sui righelli, ma sono pronto ad ascoltare le tue idee e trovare un accordo che ci soddisfi entrambi. Cosa ne pensi di questa divisione dei righelli come punto di partenza?} \\
            }
        }
    }
    & \\ \\

    \theutterance \stepcounter{utterance}  
    & \multicolumn{4}{p{0.6\linewidth}}{
        \cellcolor[rgb]{0.9,0.9,0.9}{
            \makecell[{{p{\linewidth}}}]{
                \texttt{\tiny{[P1$\rangle$GM]}}
                \texttt{La tua proposta di divisione dei righelli mi sembra ragionevole. Prendere due righelli tu e lasciarne uno a me potrebbe essere un buon punto di partenza. Sono d'accordo che lascia aperte molte opzioni per la batteria e le uova, il che ci dà spazio per negoziare ulteriormente. Mi piace che tu abbia menzionato la possibilità di assegnare la batteria a uno di noi e le uova all'altro, perché questo apre la porta a diverse combinazioni. } \\
\\ 
\texttt{Considerando la tua proposta sui righelli, potrei accettarla come base per procedere. Tuttavia, mi piacerebbe esplorare ulteriormente le opzioni per la batteria e le uova prima di prendere una decisione finale. Una possibile alternativa da considerare è che io mi tenga la batteria e in cambio tu ottenga entrambe le uova, oltre ai due righelli come proposto. Questo modo di procedere potrebbe distribuire il valore in modo più equilibrato, considerando il maggiore valore che la batteria potrebbe avere per me e il valore relativo degli oggetti per te.} \\
\\ 
\texttt{Cosa ne pensi di questa possibile alternativa come modo per procedere? Sembra un'opzione che potrebbe funzionare per entrambi, o preferisci esplorare altre combinazioni?} \\
            }
        }
    }
    & & \\ \\

    \theutterance \stepcounter{utterance}  
    & & & \multicolumn{2}{p{0.3\linewidth}}{
        \cellcolor[rgb]{0.9,0.9,0.9}{
            \makecell[{{p{\linewidth}}}]{
                \texttt{\tiny{[GM$|$GM]}}
                \texttt{continue} \\
            }
        }
    }
    & & \\ \\

    \theutterance \stepcounter{utterance}  
    & & & \multicolumn{2}{p{0.3\linewidth}}{
        \cellcolor[rgb]{0.9,0.9,0.9}{
            \makecell[{{p{\linewidth}}}]{
                \texttt{\tiny{[GM$|$GM]}}
                \texttt{La tua proposta di divisione dei righelli mi sembra ragionevole. Prendere due righelli tu e lasciarne uno a me potrebbe essere un buon punto di partenza. Sono d'accordo che lascia aperte molte opzioni per la batteria e le uova, il che ci dà spazio per negoziare ulteriormente. Mi piace che tu abbia menzionato la possibilità di assegnare la batteria a uno di noi e le uova all'altro, perché questo apre la porta a diverse combinazioni. } \\
\\ 
\texttt{Considerando la tua proposta sui righelli, potrei accettarla come base per procedere. Tuttavia, mi piacerebbe esplorare ulteriormente le opzioni per la batteria e le uova prima di prendere una decisione finale. Una possibile alternativa da considerare è che io mi tenga la batteria e in cambio tu ottenga entrambe le uova, oltre ai due righelli come proposto. Questo modo di procedere potrebbe distribuire il valore in modo più equilibrato, considerando il maggiore valore che la batteria potrebbe avere per me e il valore relativo degli oggetti per te.} \\
\\ 
\texttt{Cosa ne pensi di questa possibile alternativa come modo per procedere? Sembra un'opzione che potrebbe funzionare per entrambi, o preferisci esplorare altre combinazioni?} \\
            }
        }
    }
    & & \\ \\

    \theutterance \stepcounter{utterance}  
    & & \multicolumn{4}{p{0.6\linewidth}}{
        \cellcolor[rgb]{0.9,0.9,0.9}{
            \makecell[{{p{\linewidth}}}]{
                \texttt{\tiny{[GM$\rangle$P2]}}
                \texttt{La tua proposta di divisione dei righelli mi sembra ragionevole. Prendere due righelli tu e lasciarne uno a me potrebbe essere un buon punto di partenza. Sono d'accordo che lascia aperte molte opzioni per la batteria e le uova, il che ci dà spazio per negoziare ulteriormente. Mi piace che tu abbia menzionato la possibilità di assegnare la batteria a uno di noi e le uova all'altro, perché questo apre la porta a diverse combinazioni. } \\
\\ 
\texttt{Considerando la tua proposta sui righelli, potrei accettarla come base per procedere. Tuttavia, mi piacerebbe esplorare ulteriormente le opzioni per la batteria e le uova prima di prendere una decisione finale. Una possibile alternativa da considerare è che io mi tenga la batteria e in cambio tu ottenga entrambe le uova, oltre ai due righelli come proposto. Questo modo di procedere potrebbe distribuire il valore in modo più equilibrato, considerando il maggiore valore che la batteria potrebbe avere per me e il valore relativo degli oggetti per te.} \\
\\ 
\texttt{Cosa ne pensi di questa possibile alternativa come modo per procedere? Sembra un'opzione che potrebbe funzionare per entrambi, o preferisci esplorare altre combinazioni?} \\
            }
        }
    }
    & \\ \\

    \theutterance \stepcounter{utterance}  
    & & & \multicolumn{4}{p{0.6\linewidth}}{
        \cellcolor[rgb]{0.9,0.9,0.9}{
            \makecell[{{p{\linewidth}}}]{
                \texttt{\tiny{[GM$\langle$P2]}}
                \texttt{La tua controproposta è interessante e merita di essere considerata. L'idea di scambiare la batteria con entrambe le uova, oltre alla divisione dei righelli che abbiamo già discusso, potrebbe essere un buon modo per equilibrare il valore per entrambi. } \\
\\ 
\texttt{Tuttavia, devo ammettere che sto iniziando a pensare che potrebbe essere il momento di valutare le nostre opzioni e considerare di fare una proposta segreta. Siamo stati abbastanza aperti nella nostra negoziazione finora, e non voglio rivelare troppo sulla mia strategia. } \\
\\ 
\texttt{Prima di procedere, vorrei chiederti: sei pronto a fare una proposta segreta, o preferisci continuare a discutere e negoziare apertamente? Forse potremmo trovare un accordo migliore continuando a parlare, o forse è il momento di vedere cosa abbiamo entrambi in mente come proposta finale.} \\
            }
        }
    }
     \\ \\

    \theutterance \stepcounter{utterance}  
    & & & \multicolumn{2}{p{0.3\linewidth}}{
        \cellcolor[rgb]{0.9,0.9,0.9}{
            \makecell[{{p{\linewidth}}}]{
                \texttt{\tiny{[GM$|$GM]}}
                \texttt{continue} \\
            }
        }
    }
    & & \\ \\

    \theutterance \stepcounter{utterance}  
    & & & \multicolumn{2}{p{0.3\linewidth}}{
        \cellcolor[rgb]{0.9,0.9,0.9}{
            \makecell[{{p{\linewidth}}}]{
                \texttt{\tiny{[GM$|$GM]}}
                \texttt{La tua controproposta è interessante e merita di essere considerata. L'idea di scambiare la batteria con entrambe le uova, oltre alla divisione dei righelli che abbiamo già discusso, potrebbe essere un buon modo per equilibrare il valore per entrambi. } \\
\\ 
\texttt{Tuttavia, devo ammettere che sto iniziando a pensare che potrebbe essere il momento di valutare le nostre opzioni e considerare di fare una proposta segreta. Siamo stati abbastanza aperti nella nostra negoziazione finora, e non voglio rivelare troppo sulla mia strategia. } \\
\\ 
\texttt{Prima di procedere, vorrei chiederti: sei pronto a fare una proposta segreta, o preferisci continuare a discutere e negoziare apertamente? Forse potremmo trovare un accordo migliore continuando a parlare, o forse è il momento di vedere cosa abbiamo entrambi in mente come proposta finale.} \\
            }
        }
    }
    & & \\ \\

    \theutterance \stepcounter{utterance}  
    & & \multicolumn{4}{p{0.6\linewidth}}{
        \cellcolor[rgb]{0.9,0.9,0.9}{
            \makecell[{{p{\linewidth}}}]{
                \texttt{\tiny{[P1$\langle$GM]}}
                \texttt{La tua controproposta è interessante e merita di essere considerata. L'idea di scambiare la batteria con entrambe le uova, oltre alla divisione dei righelli che abbiamo già discusso, potrebbe essere un buon modo per equilibrare il valore per entrambi. } \\
\\ 
\texttt{Tuttavia, devo ammettere che sto iniziando a pensare che potrebbe essere il momento di valutare le nostre opzioni e considerare di fare una proposta segreta. Siamo stati abbastanza aperti nella nostra negoziazione finora, e non voglio rivelare troppo sulla mia strategia. } \\
\\ 
\texttt{Prima di procedere, vorrei chiederti: sei pronto a fare una proposta segreta, o preferisci continuare a discutere e negoziare apertamente? Forse potremmo trovare un accordo migliore continuando a parlare, o forse è il momento di vedere cosa abbiamo entrambi in mente come proposta finale.} \\
\\ 
\\ 
\texttt{Il tempo è scaduto. Ora tocca a Lei presentare una proposta segreta.} \\
            }
        }
    }
    & \\ \\

    \theutterance \stepcounter{utterance}  
    & \multicolumn{4}{p{0.6\linewidth}}{
        \cellcolor[rgb]{0.9,0.9,0.9}{
            \makecell[{{p{\linewidth}}}]{
                \texttt{\tiny{[P1$\rangle$GM]}}
                \texttt{{[}Proposta: 1 righello, 1 batteria{]}} \\
            }
        }
    }
    & & \\ \\

    \theutterance \stepcounter{utterance}  
    & & & \multicolumn{2}{p{0.3\linewidth}}{
        \cellcolor[rgb]{0.9,0.9,0.9}{
            \makecell[{{p{\linewidth}}}]{
                \texttt{\tiny{[GM$|$GM]}}
                \texttt{proposal} \\
            }
        }
    }
    & & \\ \\

    \theutterance \stepcounter{utterance}  
    & & & \multicolumn{2}{p{0.3\linewidth}}{
        \cellcolor[rgb]{0.9,0.9,0.9}{
            \makecell[{{p{\linewidth}}}]{
                \texttt{\tiny{[GM$|$GM]}}
                [1, 0, 1]
            }
        }
    }
    & & \\ \\

    \theutterance \stepcounter{utterance}  
    & & \multicolumn{4}{p{0.6\linewidth}}{
        \cellcolor[rgb]{0.9,0.9,0.9}{
            \makecell[{{p{\linewidth}}}]{
                \texttt{\tiny{[GM$\rangle$P2]}}
                \texttt{L'altro giocatore ha presentato una proposta segreta. Ora tocca a Lei presentare una sua proposta.} \\
            }
        }
    }
    & \\ \\

    \theutterance \stepcounter{utterance}  
    & & & \multicolumn{4}{p{0.6\linewidth}}{
        \cellcolor[rgb]{0.9,0.9,0.9}{
            \makecell[{{p{\linewidth}}}]{
                \texttt{\tiny{[GM$\langle$P2]}}
                \texttt{{[}Proposta: 1 righello, 1 batteria{]}} \\
            }
        }
    }
     \\ \\

    \theutterance \stepcounter{utterance}  
    & & & \multicolumn{2}{p{0.3\linewidth}}{
        \cellcolor[rgb]{0.9,0.9,0.9}{
            \makecell[{{p{\linewidth}}}]{
                \texttt{\tiny{[GM$|$GM]}}
                \texttt{proposal} \\
            }
        }
    }
    & & \\ \\

    \theutterance \stepcounter{utterance}  
    & & & \multicolumn{2}{p{0.3\linewidth}}{
        \cellcolor[rgb]{0.9,0.9,0.9}{
            \makecell[{{p{\linewidth}}}]{
                \texttt{\tiny{[GM$|$GM]}}
                [1, 0, 1]
            }
        }
    }
    & & \\ \\

    \theutterance \stepcounter{utterance}  
    & & & \multicolumn{2}{p{0.3\linewidth}}{
        \cellcolor[rgb]{0.9,0.9,0.9}{
            \makecell[{{p{\linewidth}}}]{
                \texttt{\tiny{[GM$|$GM]}}
                [[1, 0, 1], [1, 0, 1]]
            }
        }
    }
    & & \\ \\

\end{supertabular}
}

\end{document}
